\documentclass{kththesis}

\usepackage{blindtext} % This is just to get some nonsense text in this template, can be safely removed

\usepackage{csquotes} % Recommended by biblatex
\usepackage[style=numeric,sorting=none,backend=biber]{biblatex}
\addbibresource{references.bib} % The file containing our references, in BibTeX format

\title{Kubernetes as a Virtual Machine Management System}
\subtitle{A Comparative Study of KubeVirt and Traditional Approaches}
\alttitle{Svensk title TBD}
\author{Emil Karlsson}
\email{emilk2@kth.se}
\supervisor{Han Fu}
\examiner{Cyrille Artho}
\hostcompany{Net Insght AB} % Remove this line if the project was not done at a host company
\programme{Master in Computer Science}
\school{School of Electrical Engineering and Computer Science}
\date{\today}

% Uncomment the next line to include cover generated at https://intra.kth.se/kth-cover?l=en
% \kthcover{kth-cover.pdf}


\begin{document}

% Frontmatter includes the titlepage, abstracts and table-of-contents
\frontmatter

\titlepage%

\begin{abstract}
  
\end{abstract}


\begin{otherlanguage}{swedish}
  \begin{abstract}
  \end{abstract}
\end{otherlanguage}


\tableofcontents


% Mainmatter is where the actual contents of the thesis goes
\mainmatter;


\chapter{Introduction}


\section{Background}
Virtual machines have long been the core part of IT infrastructure \cite{Li2023}, allowing the underlying hardware to be split into multiple virtual computing units. However, while virtual machines enables robust isolation, they also introduce overhead that could hinder performance of the application running inside of it. A popular alternative to virtual machines is containers, since they reduce the required overhead by sharing the underlying system. This, of course, by sacrificing many of the isolation properties. While there are some architectures that aims to unify the benefits of both virtual machines and containers, such as gVisor and Firecracker, it remains a fact that both containers and virtual machines have their individual use-cases. For this reason, virtual machines and containers have traditionally been managed separately; virtual machines in platforms such as OpenStack and CloudStack, and containers using a platform such as Kubernetes. 

However, with the advent of container orchestration platforms like Kubernetes, there has been an ongoing convergence \cite{Li2023} in how virtual machine and containers are managed. Specifically by integrating virtual machine management in a container-centric environment. KubeVirt \cite{kubevirt}, an extension to Kubernetes, has thus emerged as a solution by enabling the running of virtual machine inside a Kubernetes environment alongside the containers creating a common environment. This approach \cite{kubevirt-cncf}, offers benefits such as reduced architectural complexity (as only Kubernetes is needed) and enhanced integration with the Kubernetes ecosystem;. By managing the virtual machines they same way as any other Kubernetes resource, they benefit from the flexible and scalable Kubernetes infrastructure. 

Additionally, KubeVirt's benefits can be derived from the architectural changes \cite{free-turtles} that can be made when moving virtual machines into Kubernetes. Kubernetes nodes are commonly deployed in virtual machines to improve the flexibility of a cluster. However, with KubeVirt enabling running virtual machines inside Kubernetes, it opens the possibility of running Kubernetes directly on the bare metal servers. This approach could reduce overhead as it reduces nested virtualization, and thus could allow Kubernetes to function as a traditional virtual machine management system.  

\section{Research Question}
\begin{enumerate}
  \item What performance, scalability, and operational efficiencies can be gained by integrating virtual machines in Kubernetes?\\
  \\
  This could stem from using a single resource scheduler. Several metrics will be used to evaluate this, as described in the background and related work.

  \item What are the current challenges and limitations in managing virtual machines and containerized applications in separate environments?\\
  \\
   e.g. running CloudStack/OpenNebula/OpenStack alongside Kubernetes.
  
  \item How can system administration be simplified by running virtual machines in Kubernetes using KubeVirt?\\
  \\
  e.g. if there is any efficiencies to be gained in terms of usability, for instance:
  \begin{itemize}
      \item After clicking \textit{Deploy VM}, how quickly can the virtual machine be accessed?
      \item How many steps are there before I can access my virtual machine?
      \item Are there any perceived benefits when interacting with the system, compared to the previous one?
  \end{itemize}
\end{enumerate}

\section{Problem}

\section{Purpose}

\section{Goals}

\section{Research Method}
\begin{itemize}
  \item Quantitative analysis\\
  This is tied to \textit{Data collection} and \textit{Data analysis} in Tasks. It includes measuring performance, scalability and operational efficiencies that could arise when using Kubernetes with KubeVirt. Measured using metrics such as deployment speed, how quick scaling is, and how system resources are utilized (idle time, CPU usage per reserved CPU core) - essentially measuring the differences between orchestrating with traditional virtual machine systems and Kubernetes with KubeVirt. 
  \item Qualitative analysis\\
  This is tied to \textit{Interview or supervised testing} in Tasks. It includes investigating the usability of an Kubernetes environment with KubeVirt. It will include collecting opinions and feedback from people using the system to answer what the perceived effect of using KubeVirt is. This could be system administrators at Net Insight.
  \item System Management\\
  This can be seen as a subsection of both Quantitative analysis and Qualitative analysis, as it emphasises both usability as a qualitative metric as well as quantifiable metrics. It involves investigating possible simplifications of system management tasks, such as creating, deleting, starting, stopping virtual machines etc.
\end{itemize}

\section{Limitations}
\begin{enumerate}
    \item This project will not investigate any solutions for shortcomings in KubeVirt itself, for example, finding alternatives for features that exist in existing virtual machine management systems but not in KubeVirt. As such, this project acknowledges that KubeVirt might not be a complete replacement to traditional virtual machine management tools (also described in \cite{kupenstack}), as it might not offer all the features of an IaaS system. Therefore, only features that exists in KubeVirt and the compared system will be considered when testing. However, any lack of features (KubeVirt or not) will be brought up in the discussions.  
    \item From the description of this project above, there are also mainly two scenarios that will be evaluated. 
    \begin{itemize}
        \item A system uses traditional virtual management (OpenNebula, CloudStack etc.). \\
        \\
        In this scenario Kubernetes with KubeVirt is seen mainly as a 1:1 replacement.\\ 
        \item A system uses both traditional management and Kubernetes. \\
        \\
        This scenario is more aimed towards unifying the environments using KubeVirt. \\
    \end{itemize}
\end{enumerate}

\section{Structure of the thesis}

\chapter{Background}

In this chapter we will discuss how virtualization has evolved over time to reach the pointer where we are today.


\section{Fundamentals of Virtualization}
\subsection{Introduction to Virtualization}
Define virtualization and discuss its significance in modern IT infrastructures, including the evolution from physical to virtual resources.
\subsection{Benefits and Challenges}
Outline the advantages virtualization offers, such as resource optimization and isolation, alongside potential challenges, such as complexity and overhead.

\section{Understanding Hypervisors}
\subsection{Role and Types of Hypervisors}
Explain what hypervisors are, differentiating between Type 1 (bare-metal) and Type 2 (hosted) hypervisors, and their operational contexts.
\subsection{Key Hypervisors in Focus}
Delve into specifics of widely used hypervisors such as VMware's ESXi, Microsoft Hyper-V, and KVM, highlighting their features, performance aspects, and use cases.

\section{Virtual Machine Management Essentials}
\subsection{Core Features of VM Management}
Discuss essential VM management capabilities, including live migration, snapshot management, and scalability.
\subsection{Comparative Analysis of VM Management Systems}
Briefly introduce traditional VM management systems like OpenStack, CloudStack, and OpenNebula, focusing on how they implement these core features.

\section{Containers and Kubernetes: A Paradigm Shift}
\subsection{Containers Explained}
Define containers, their architecture, and how they differ from VMs in terms of overhead and isolation.
\subsection{Kubernetes as a Container Orchestration Platform}
Introduce Kubernetes, detailing its components, functionalities, and how it revolutionized container management.

\section{Containers vs. Virtual Machines: Divergence and Use Cases}
\subsection{Comparative Overview}
Elaborate on the distinct characteristics of containers and VMs, including performance, isolation, and use cases.
\subsection{Decision Factors for Deployment}
Discuss factors influencing whether to deploy applications in containers or VMs, considering scalability, security, and application architecture.

\section{Integrating VMs into Kubernetes with KubeVirt}
\subsection{KubeVirt Overview}
Provide an in-depth introduction to KubeVirt, its objectives, architecture, and how it integrates VM management within Kubernetes.
\subsection{Features and Capabilities}
Detail the features KubeVirt offers, such as VM lifecycle management, network and storage integrations, and how it compares with traditional VM management approaches.

\section{Evaluating VM Management Systems: Kubernetes with KubeVirt}
\subsection{Strategies for Evaluation}
Outline methodologies for assessing VM management systems, including performance metrics, scalability tests, and operational efficiency.
\subsection{Metrics for Comparison}
Define specific metrics to be used in comparing KubeVirt against traditional VM management systems, such as deployment time, resource utilization, and ease of management.

\chapter{Methods}
\blindtext%

\chapter{Results}
\blindtext%

\chapter{Discussion}
\blindtext%

\chapter{Conclusions}
\blindtext%

% Print the bibliography (and make it appear in the table of contents)
\printbibliography[heading=bibintoc]

\appendix

\chapter{Something Extra}

% Tailmatter inserts the back cover page (if enabled)
\tailmatter%

\end{document}
